\documentclass[10pt,a4paper]{amsart}
\usepackage[margin=19mm]{geometry}
\usepackage{amsmath,amssymb,mathtools}
\usepackage[scr=rsfs,cal=euler]{mathalfa}
\usepackage{xfrac}
\usepackage[unicode,hidelinks,bookmarks=false]{hyperref}
\newcommand{\ie}{i.e.\ }
\newcommand{\cf}{cf.\ }
\newcommand{\resp}{respectively\ }
\newcommand{\Subsection}{\subsection}
\providecommand{\nobreakdash}{\nobreak\hbox{-}}
\setlength{\emergencystretch}{2em}
\multlinegap0pt
\usepackage{bm}
\usepackage{bbm}
\usepackage{dsfont}
\usepackage{xspace}
\usepackage{cancel}
\usepackage{mathtools}
\usepackage{amsthm}
\makeatletter
\@ifundefined{theorem}{\newtheorem{theorem}{Theorem}}{}
\makeatother
\title[On the Gromov-Hausdorff distance between the cloud of bounde]{On the Gromov-Hausdorff distance between the cloud of bounded metric spaces and a cloud with nontrivial stabilizer}
\author{B. A. Nesterov}
\date{Source: arXiv:2505.23563v1 (CC BY 4.0)}
\begin{document}
\maketitle

\noindent\emph{Source and licence.} On the Gromov-Hausdorff distance between the cloud of bounded metric spaces and a cloud with nontrivial stabilizer. Version arXiv:2505.23563v1 (CC BY 4.0), distributed under the Creative Commons Attribution 4.0 International licence, as recorded in the arXiv licence metadata for this exact version and shown on the abstract page https://arxiv.org/abs/2505.23563v1. Full rights notice: licenses/arxiv-2505.23563-CC-BY-4.0.txt.\par\smallskip

\noindent\textit{Editorial scope.} Counted unit: the Theorem whose source label is \texttt{thrmCenterImage} in the article (source0.tex lines 242--244, source label \texttt{thrmCenterImage}), delivered together with the article's own complete original proof of it (source0.tex lines 246--302, one \texttt{proof} environment of 57 lines). No mathematics is written, repaired or re-derived: statement and proof are the article's own text line for line. Required context retained uncounted: none. Every definition, display and label of those lines is retained unchanged. Omitted from this package and not redistributed: 1--241, 245--245, 303--475. Nothing inside a retained excerpt is omitted or altered: every retained body is delivered exactly as the article's lines are written, so no omission receipt of any kind accompanies this package. Citations: the retained excerpts carry no citation command at all, so no bibliography entry of the article is transcribed and no reference is redistributed. Reference adaptation: none was needed and none was made; every reference inside the delivered unit resolves inside it, so no unresolved reference or citation is produced. Printed slips kept literal: no printed word, symbol, hypothesis or number of the counted statement or of its proof was altered. Layout and class: the article's own class is \texttt{article} with packages including geometry, graphicx, inputenc, babel, mathtools, amsfonts, amssymb, mathrsfs, amscd, amsmath, amsthm, verbatim, url, stmaryrd; the wrapper is the lane's pinned \texttt{amsart} editorial wrapper at 10pt on A4, which restates the theorem-like environments the retained text needs and adds the article's own macro definitions that the retained text needs (none), with extra wrapper packages (bm, bbm, dsfont, xspace, cancel, mathtools). The delivered numbering is the unit's own. Assets: no retained span contains a figure environment, a graphics-inclusion command, a PDF-page inclusion, a table or a TikZ picture; no image file is copied into this package, the package declares no asset dependency and no image file is redistributed with it. Licence: version arXiv:2505.23563v1 of this article is distributed under the Creative Commons Attribution 4.0 International licence, as recorded in the arXiv licence metadata for this exact version and shown on the abstract page https://arxiv.org/abs/2505.23563v1. A full rights notice is delivered at licenses/arxiv-2505.23563-CC-BY-4.0.txt. Characters: every retained excerpt is pure ASCII, so the delivered engine, XeLaTeX with its default Latin Modern text font, reports no missing character.
\begin{theorem}\label{thrmCenterImage}
Let $M$ be the center of the cloud $[M]$ with a nontrivial stabilizer. Let $R$ be a correspondence between $[\Delta_{1}]$ and $[M]$ with finite distortion $\varepsilon$. Then any space from the image $R(\Delta_{1})$ lies at a distance from $M$ not exceeding $2\varepsilon$.
\end{theorem}

\begin{proof}
The nontriviality of the stabilizer $[M]$ implies that there exists a number $l > 1$ such that $\{l^{j} \mid j\in \mathbb{Z}\}$ is a subgroup of $\operatorname{St}([M])$.

Fix $Y$ in the image of $\Delta_{1}$. Suppose that $|M, Y| = d > \varepsilon$. Denote $|Y, kY| = \rho$, where $k \geq 2$, $k = l^{j_{1}}$. By the triangle inequality $\rho + d \geq kd$, hence $\rho \geq (k-1)d > (k-1)\varepsilon$. Then $kY$ lies in the image of $X \neq \Delta_1$, with $\rho - \varepsilon \leq |X, \Delta_1| \leq \rho + \varepsilon$.

Take arbitrary $\alpha > 0$ and $\beta \in (0,1)$. For spaces $(1+\alpha)X$, $(1-\beta)X$ the following inequalities hold:
\begin{align*}
|X, (1+\alpha)X| &= \alpha |X, \Delta_1| \leq \alpha\rho + \alpha\varepsilon, \\
|X, (1-\beta)X| &= \beta|X, \Delta_1| \leq \beta\rho + \beta\varepsilon, \\
|(1+\alpha) X, (1-\beta)X| &= (\alpha + \beta)|X, \Delta_1| \geq (\alpha+\beta)\rho - (\alpha+\beta)\varepsilon.
\end{align*}

There exist $Y_\alpha, Y_\beta \in [M]$ such that $kY_\alpha \in R\big((1+\alpha)X\big)$, $kY_\beta \in R\big((1-\beta)X\big)$, and for them the following inequalities hold:
\begin{align*}
|kY, kY_\alpha| &\leq |X, (1+\alpha)X| + \varepsilon \leq \alpha\rho + (\alpha+1)\varepsilon, \\
|kY, kY_\beta| &\leq |X, (1-\beta)X| + \varepsilon \leq \beta\rho + (\beta+1)\varepsilon, \\
|kY_\alpha, kY_\beta| &\geq |(1+\alpha)X, (1-\beta)X| - \varepsilon \geq (\alpha+\beta)\rho - (\alpha+\beta+1)\varepsilon.
\end{align*}

Dividing these inequalities by $k$, we obtain:
\begin{align*}
|Y, Y_{\alpha}| &\leq \frac{\alpha}{k}\rho + \frac{\alpha+1}{k}\varepsilon, \\
|Y, Y_{\beta}| &\leq \frac{\beta}{k}\rho + \frac{\beta+1}{k}\varepsilon, \\
|Y_\alpha, Y_{\beta}| &\geq \frac{\alpha+\beta}{k}\rho - \frac{\alpha+\beta+1}{k}\varepsilon.
\end{align*}

Taking preimages of spaces $Y, Y_{\alpha}, Y_{\beta}$, we obtain:
\begin{align*}
|\Delta_1, X_{\alpha}| &\leq \frac{\alpha}{k}\rho + \left(\frac{\alpha+1}{k} + 1\right)\varepsilon, \\
|\Delta, X_{\beta}| &\leq \frac{\beta}{k}\rho + \left(\frac{\beta+1}{k}+1\right)\varepsilon, \\
|X_\alpha, X_{\beta}| &\geq \frac{\alpha+\beta}{k}\rho - \left(\frac{\alpha+\beta+1}{k}+1\right)\varepsilon.
\end{align*}

Assuming $\alpha > \beta$ we obtain:
$$
\frac{\alpha+\beta}{k}\rho - \left(\frac{\alpha+\beta+1}{k}+1\right)\varepsilon \leq \frac{\alpha}{k}\rho + \left(\frac{\alpha+1}{k} + 1\right)\varepsilon,
$$
which is equivalent to:
$$
\rho \leq \frac{k}{\beta}\left(\frac{2\alpha+\beta+2}{k}+2\right)\varepsilon,
$$
and further:
$$
\rho \leq \left(1+\frac{2\alpha + 2}{\beta} + 2\frac{k}{\beta}\right)\varepsilon.
$$

We are interested in an upper bound for $d$:
$$
d \leq \frac{\rho}{k-1} \leq \left(\frac{1}{k-1}+\frac{2\alpha + 2}{\beta(k-1)} + 2\frac{k}{\beta(k-1)}\right)\varepsilon.
$$

The last term in parentheses is strictly greater than 2 for any $k>2$, $\alpha>0$, $\beta\in (0,1)$, while the other terms tend to 0 as $k$ increases. Since the stabilizer is nontrivial, it contains sequences of numbers tending to both 0 and $\infty$. Taking $\beta$ to 1 and $k$ to infinity, we obtain the estimate:
$$
|Y, M| \leq 2\varepsilon,
$$
which completes the proof.
\end{proof}

\end{document}
